\Needspace{9\baselineskip}
\begin{center}\textit{Solution of the equation \(axx+byy+czz=0\).}\end{center}
\Article{294}
\textsc{Theorem.} \textit{Let \(a,b,c\) be pairwise coprime numbers, none
of which is \(=0\) or divisible by a square. The equation}
\[
axx+byy+czz=0\ldots(\Omega)
\]
\textit{will have no integer solution (apart from \(x=y=z=0\), which we
disregard) unless \(-bc,-ac,-ab\) are quadratic residues of \(a,b,c\), respectively,
and these numbers do not all have the same sign. But when these four
conditions hold, \((\Omega)\) will be solvable in integers.}

\textit{Proof.} If \((\Omega)\) is solvable in integers at all, it can also
be solved by values of \(x,y,z\) having no common divisor. For any values
satisfying the equation \(\Omega\) still satisfy it when divided by their
greatest common divisor. Now suppose \(app+bqq+crr=0\) and that \(p,q,r\) have
no common divisor. They will then also be pairwise coprime. For if
\(q,r\) had a common divisor \(\mu\), this would be coprime to \(p\),
but \(\mu\mu\) would divide \(app\), and therefore also \(a\), contrary
to the hypothesis. Likewise, \(p,r;\ p,q\) will be pairwise coprime. So
\(-app\) is represented by the binary form \(byy+czz\),
\SourceLeaf{355}
on assigning to \(y,z\) the coprime values \(q,r\). Its determinant
\(-bc\) will therefore be a quadratic residue of \(app\), and hence
also of \(a\) (Article 154). In the same way, \(-acRb,-abRc\). That \((\Omega)\)
cannot have a solution if \(a,b,c\) all have the same sign is so obvious
that it needs no explanation.

We will arrange the proof of the converse, which forms the second
part of the theorem, as follows. \textit{First}, we show how to find
a ternary form equivalent to \(\left(\begin{smallmatrix}a,b,c\\0,0,0\end{smallmatrix}\right)\ldots f\) whose coefficients \(2,3,4\) are
divisible by \(abc\). \textit{Second}, from this we will derive a
solution of the equation \((\Omega)\).

I.\ Find three integers \(A,B,C\) having no common divisor and such
that \(A\) is coprime to \(b\) and \(c\); \(B\) to \(a\)
and \(c\); \(C\) to \(a\) and \(b\); while \(aAA+bBB+cCC\) is
divisible by \(abc\). This is done as follows. Let \(\mathfrak A,\mathfrak B,\mathfrak C\) be values
of the expressions \(\surd-bc\pmod a\), \(\surd-ac\pmod b\), \(\surd-ab\pmod c\), respectively. They
will necessarily be coprime to \(a,b,c\), respectively. Choose three
integers \(\mathfrak a,\mathfrak b,\mathfrak c\) completely at will, provided they are coprime to
\(a,b,c\), respectively (\textit{for example}, all \(=1\)), and determine
\(A,B,C\) so that
\[
\begin{aligned}
A&\equiv\mathfrak b c\pmod b&&\text{and }\equiv\mathfrak c\mathfrak C\pmod c\\
B&\equiv\mathfrak c a\pmod c&&\text{and }\equiv\mathfrak a\mathfrak A\pmod a\\
C&\equiv\mathfrak a b\pmod a&&\text{and }\equiv\mathfrak b\mathfrak B\pmod b
\end{aligned}
\]
Then
\[
\begin{aligned}
aAA+bBB+cCC&\equiv\mathfrak a\mathfrak a(b\mathfrak A\mathfrak A+cbb)\\
&\equiv\mathfrak a\mathfrak a(b\mathfrak A\mathfrak A-\mathfrak A\mathfrak A b)\equiv0\pmod a
\end{aligned}
\]
so it will be divisible by \(a\), and likewise by \(b,c\), and
therefore also by \(abc\). Further, \(A\) is clearly necessarily
coprime to \(b\) and \(c\); \(B\) to \(a\) and \(c\);
\(C\) to \(a\) and \(b\). But if these values of \(A,B,C\) have
a (greatest) common divisor \(\mu\), this will clearly be coprime
to \(a,b,c\), and so to \(abc\). Dividing those values by \(\mu\)
therefore gives new ones with no common divisor, still making
\(aAA+bBB+cCC\) divisible by \(abc\), and so satisfying all the conditions.

II.\ With \(A,B,C\) determined in this way, \(Aa,Bb,Cc\) will also have
no common divisor. For if they had a common divisor \(\mu\), it
would necessarily be coprime to \(a\) (since the latter is coprime
to both \(Bb\) and \(Cc\)), and similarly to
\SourceLeaf{356}
\(b\) and \(c\). So \(\mu\) would also have to divide \(A,B,C\),
contrary to the hypothesis. We can therefore find integers \(\alpha,\beta,\gamma\)
such that \(\alpha Aa+\beta Bb+\gamma Cc=1\). Find, further, six integers \(\alpha',\beta',\gamma',\alpha'',\beta'',\gamma''\) such that
\[
\beta'\gamma''-\gamma'\beta''=Aa,\quad
\gamma'\alpha''-\alpha'\gamma''=Bb,\quad
\alpha'\beta''-\beta'\alpha''=Cc
\]
Now let \(f\) go by the substitution
\[
\begin{array}{ccc}
\alpha,&\alpha',&\alpha''\\
\beta,&\beta',&\beta''\\
\gamma,&\gamma',&\gamma''
\end{array}
\]
into \(\left(\begin{smallmatrix}m,m',m''\\n,n',n''\end{smallmatrix}\right)=g\) (which will be equivalent to \(f\)). I say that
\(m',m'',n\) will be divisible by \(abc\). For set
\[
\begin{aligned}
\beta''\gamma-\gamma''\beta&=A',&\gamma''\alpha-\alpha''\gamma&=B',&\alpha''\beta-\beta''\alpha&=C'\\
\beta\gamma'-\gamma\beta'&=A'',&\gamma\alpha'-\alpha\gamma'&=B'',&\alpha\beta'-\beta\alpha'&=C''
\end{aligned}
\]
and we have
\[
\begin{aligned}
\alpha'&=B''Cc-C''Bb,&\beta'&=C''Aa-A''Cc,&\gamma'&=A''Bb-B''Aa\\
\alpha''&=C'Bb-B'Cc,&\beta''&=A'Cc-C'Aa,&\gamma''&=B'Aa-A'Bb
\end{aligned}
\]
Substituting these values into the equations
\[
\begin{aligned}
m'&=a\alpha'\alpha'+b\beta'\beta'+c\gamma'\gamma'\\
m''&=a\alpha''\alpha''+b\beta''\beta''+c\gamma''\gamma''\\
n&=a\alpha'\alpha''+b\beta'\beta''+c\gamma'\gamma''
\end{aligned}
\]
gives, modulo \(a\),
\[
\begin{aligned}
m'&\equiv bcA''A''(BBb+CCc)\equiv0\\
m''&\equiv bcA'A'(BBb+CCc)\equiv0\\
n&\equiv bcA'A''(BBb+CCc)\equiv0
\end{aligned}
\]
that is, \(m',m'',n\) will be divisible by \(a\). In the same way,
these numbers are found to be divisible by \(b,c\), and therefore
also by \(abc\), as claimed.
\SourceLeaf{357}
\par
III.\ For convenience, write the determinant of the forms \(f,g\)
as \(-abc=d\), and set
\[
md=M,\quad m'=M'd,\quad m''=M''d,\quad n=Nd,\quad n'=N',\quad n''=N''
\]
Clearly \(f\) goes by the substitution \((S)\)
\[
\begin{array}{ccc}
\alpha d,&\alpha',&\alpha''\\
\beta d,&\beta',&\beta''\\
\gamma d,&\gamma',&\gamma''
\end{array}
\]
into the ternary form \(\left(\begin{smallmatrix}Md,M'd,M''d\\Nd,N'd,N''d\end{smallmatrix}\right)=g'\) of determinant \(d^3\), which will
therefore be contained under \(f\). I now say that this form
\(g'\) is necessarily equivalent to \(\left(\begin{smallmatrix}d,0,0\\d,0,0\end{smallmatrix}\right)=g''\). For \(\left(\begin{smallmatrix}M,M',M''\\N,N',N''\end{smallmatrix}\right)=g'''\) is
clearly a ternary form of determinant \(1\). Further, since by
the hypothesis \(a,b,c\) do not all have the same sign, \(f\)
will be an indefinite form, from which it follows readily that
\(g'\) and \(g'''\) must also be indefinite. So \(g'''\) will
be equivalent to the form \(\left(\begin{smallmatrix}1,0,0\\1,0,0\end{smallmatrix}\right)\) (Article 277), and we can find
a transformation \((S')\) of the former into the latter. But clearly
\((S')\) takes the form \(g'\) into \(g''\). Thus \(g''\), too,
will be contained under \(f\), and combining the substitutions
\((S),(S')\) gives a transformation of the form \(f\) into \(g''\).
If this is
\[
\begin{array}{ccc}
\delta,&\delta',&\delta''\\
\varepsilon,&\varepsilon',&\varepsilon''\\
\zeta,&\zeta',&\zeta''
\end{array}
\]
we clearly have two solutions of the equation \((\Omega)\), namely
\(x=\delta',y=\varepsilon',z=\zeta'\) and \(x=\delta'',y=\varepsilon'',z=\zeta''\). Neither set of values can all be \(=0\),
since necessarily
\[
\delta\varepsilon'\zeta''+\delta'\varepsilon''\zeta+\delta''\varepsilon\zeta'-\delta\varepsilon''\zeta'-\delta'\varepsilon\zeta''-\delta''\varepsilon'\zeta=d.\quad\text{As required.}
\]

\emph{Example.} Let the given equation be \(7xx-15yy+23zz=0\), which is solvable
because \(345\,R\,7,-161\,R\,15,105\,R\,23\). Here the values of \(\mathfrak A,\mathfrak B,\mathfrak C\) are \(3,7,6\);
setting \(\mathfrak a=\mathfrak b=\mathfrak c=1\) gives \(A=98,B=-39,C=-8\). From this we obtain the substitution
\[
\left\{\begin{array}{rrr}3,&5,&22\\-1,&2,&-28\\8,&25,&-7\end{array}\right\}
\]
which takes \(f\) into \(\left(\begin{smallmatrix}1520,14490,-7245\\-2415,-1246,4735\end{smallmatrix}\right)=g\). This gives
\[
(S)=\left\{\begin{array}{rrr}7245,&5,&22\\-2415,&2,&-28\\19320,&25,&-7\end{array}\right\},\quad
g'''=\left(\begin{smallmatrix}3670800,6,-3\\-1,-1246,4735\end{smallmatrix}\right).
\]
\SourceLeaf{358}
The form \(g'''\) is found to go into \(\left(\begin{smallmatrix}1,0,0\\1,0,0\end{smallmatrix}\right)\) by the substitution
\[
\left\{\begin{array}{rrr}3,&5,&1\\-2440,&-4066,&-813\\-433,&-722,&-144\end{array}\right\}\ldots(S')
\]
Combining this with \((S)\) gives \(\left\{\begin{smallmatrix}9,11.12\\-1,9,-9\\-9,4,3\end{smallmatrix}\right\}\), which takes \(f\)
into \(g''\). So we have two solutions of the given equation,
\(x=11,y=9,z=4\) and \(x=12,y=-9,z=3\). The latter becomes simpler on dividing the
values by the common divisor \(3\), giving \(x=4,y=-3,z=1\).

\Article{295}
The latter part of the theorem in the preceding article can also
be proved as follows. Find an integer \(h\) such that \(ah\equiv\mathfrak C\pmod c\)
(we use the symbols \(\mathfrak A,\mathfrak B,\mathfrak C\) in the same sense as in the preceding
article), and set \(ahh+b=ci\). It is then easy to see that \(i\)
is an integer and that the number \(-ab\) is the determinant of
the binary form \((ac,ah,i)\ldots\varphi\). This form certainly will not be positive
(since by the hypothesis \(a,b,c\) do not all have the same sign,
\(ab\) and \(ac\) cannot both be positive). Further, it will
have characteristic number \(-1\), which we prove synthetically
as follows. Determine integers \(e,e'\) such that
\[
\begin{aligned}
e&\equiv0\pmod a\text{ and }\equiv\mathfrak B\pmod b;\\
ce'&\equiv\mathfrak A\pmod a\text{ and }\equiv h\mathfrak B\pmod b
\end{aligned}
\]
Then \((e,e')\) will be a value of the expression \(\surd-(ac,ah,i)\pmod{-ab}\). For modulo
\(a\) we have
\[
\begin{aligned}
ee&\equiv0\equiv-ac,\quad ee'\equiv0\equiv-ah\\
cce'e'&\equiv\mathfrak A\mathfrak A\equiv-bc\equiv-cci\quad\text{so }e'e'\equiv-i
\end{aligned}
\]
and modulo \(b\) we have
\[
\begin{aligned}
ee&\equiv\mathfrak B\mathfrak B\equiv-ac,\quad cee'\equiv h\mathfrak B\mathfrak B\equiv-ach\quad\text{so }ee'\equiv-ah\\
cce'e'&\equiv hh\mathfrak B\mathfrak B\equiv-achh\equiv-cci\quad\text{so }e'e'\equiv-i
\end{aligned}
\]
These same three congruences, which hold for both moduli \(a,b\),
will also hold modulo \(ab\). From the theory of ternary forms
we readily conclude that \(\varphi\) can be represented by the form
\(\left(\begin{smallmatrix}-1,0,0\\1,0,0\end{smallmatrix}\right)\). So let
\[
actt+2ahtu+iuu=-(\alpha t+\beta u)^2+2(\gamma t+\delta u)(\varepsilon t+\zeta u)
\]
\SourceLeaf{359}
Multiplying by \(c\) gives
\[
a(ct+hu)^2+buu=-c(\alpha t+\beta u)^2+2c(\gamma t+\delta u)(\varepsilon t+\zeta u)
\]
It follows that, if \(t,u\) are assigned definite values making
either \(\gamma t+\delta u\) or \(\varepsilon t+\zeta u\) become \(=0\), we have a solution
of the equation \(\Omega\). It is therefore satisfied both by
\[
x=\delta c-\gamma h,\quad y=\gamma,\quad z=\alpha\delta-\beta\gamma
\]
and by
\[
x=\zeta c-\varepsilon h,\quad y=\varepsilon,\quad z=\alpha\zeta-\beta\varepsilon
\]
Neither set of values can all become \(=0\). For if \(\delta c-\gamma h=0,\gamma=0\),
we would also have \(\delta=0\) and \(\varphi=-(\alpha t+\beta u)^2\), giving \(ab=0\), contrary
to the hypothesis; and likewise for the other set. --- In our example
we find the form \(\varphi\) to be \((161,-63,24)\), the value of the expression
\(\surd-\varphi\pmod{105}=(7,-51)\), and the following representation of the form \(\varphi\) by
\(\left(\begin{smallmatrix}-1,0,0\\1,0,0\end{smallmatrix}\right)\):
\[
\varphi=-(13t-4u)^2+2(11t-4u)(15t-5u)
\]
This gives the solutions \(x=7,y=11,z=-8\); \(x=20,y=15,z=-5\), or, on dividing by
\(5\) and disregarding the sign of \(z\), \(x=4,y=3,z=1\).

Of these two methods for solving the equation \(\Omega\), the latter
has the advantage of usually working with smaller numbers. But the
former, which can also be shortened by various devices we pass over
here, seems more elegant, especially because the numbers \(a,b,c\)
are treated in exactly the same way and the calculation is unchanged
by any permutation of them. This is different in the second method,
where the calculation is usually most convenient if we take the
smallest of the three given numbers for \(a\) and the largest
for \(c\), as we did in our example.

\begin{center}\itshape On the method by which the distinguished Le Gendre treated the fundamental theorem.\end{center}
\Article{296}
The elegant theorem discussed in the preceding articles was first
discovered by the distinguished Le Gendre, \emph{Hist. de l'Ac. de
Paris} 1784, p.\ 507, and supplied with a beautiful proof, entirely
different from our two. At the same place this excellent geometer
also tried to derive from it a proof of propositions agreeing with
the fundamental theorem of
\SourceLeaf{360}
the preceding section. We have already said above, in Article 151,
that this proof seems to us unsuitable for that purpose. This is
therefore the place to set out this proof, which is very elegant
in itself, briefly and to give the reasons for our judgement. We
begin with the following observation: \emph{If the numbers \(a,b,c\)
are all \(\equiv1\pmod4\), the equation \(axx+byy+czz=0\ldots(\Omega)\) cannot be solvable.} For it
is very easy to see that, in this case, the value of \(axx+byy+czz\) must
be either \(\equiv1\), \(\equiv2\) or \(\equiv3\pmod4\), unless \(x,y,z\) are all
taken even. So, if \(\Omega\) were solvable, it could only be by even
values of \(x,y,z\), \emph{which is absurd}: any values satisfying
the equation \(\Omega\) still satisfy it when divided by their greatest
common divisor, which must give at least one odd value. The different
cases of the theorem to be proved now reduce to the following points.

I.\ Let \(p,q\) be primes of the form \(4n+3\) (positive and unequal).
We cannot have \(p\,R\,q,q\,R\,p\) simultaneously. For if this were possible,
setting \(1=a,-p=b,-q=c\) would clearly satisfy all conditions for the
solvability of the equation \(axx+byy+czz=0\) (Article 294). But by the preceding
observation it has no solution, so the supposition cannot hold.
Proposition 7 of Article 131 follows immediately.

II.\ If \(p\) is a prime of the form \(4n+1\) and \(q\) a prime
of the form \(4n+3\), we cannot have \(q\,R\,p,p\,N\,q\) simultaneously. Otherwise
we would have \(-p\,R\,q\), and the equation \(xx+pyy-qzz=0\) would be solvable,
whereas by the preceding observation it has no solution. Cases 4
and 5 of Article 131 follow from this.

III.\ If \(p,q\) are primes of the form \(4n+1\), we cannot have
\(p\,R\,q,q\,N\,p\) simultaneously. Take another prime \(r\) of the form
\(4n+3\) that is a residue of \(q\) and of which \(p\) is a
nonresidue. Then the cases just proved (II) give \(q\,R\,r,r\,N\,p\). So, if
\(p\,R\,q,q\,N\,p\), we would have \(qr\,R\,p,pr\,R\,q,pq\,N\,r\), and therefore \(-pq\,R\,r\). Thus the
equation \(pxx+qyy-rzz=0\) would be solvable, contrary to the preceding
observation; so the supposition cannot hold. Cases 1 and 2 of
Article 131 follow from this.

This case is handled more neatly as follows. Let \(r\) be a
prime of the form \(4n+3\) of which \(p\) is a nonresidue. Then
we also have \(r\,N\,p\), and so (assuming \(p\,R\,q,q\,N\,p\)) \(qr\,R\,p\). Further,
\(-p\,R\,q,-p\,R\,r\), and therefore also \(-p\,R\,qr\). So the equation \(xx+pyy-qrzz=0\)
would be solvable, contrary to the preceding observation. From
this, and so on.

IV.\ If \(p\) is a prime of the form \(4n+1\) and \(q\) a prime
of the form \(4n+3\),
\SourceLeaf{361}
we cannot have \(p\,R\,q,q\,N\,p\) simultaneously. Take an auxiliary prime
\(r\) of the form \(4n+1\) that is a nonresidue of both \(p,q\).
Then (by II) \(q\,N\,r\), and (by III) \(p\,N\,r\); hence \(pq\,R\,r\). So,
if \(p\,R\,q,q\,N\,p\), we would also have \(pr\,N\,q,-pr\,R\,q,qr\,R\,p\), and the equation \(pxx-qyy+rzz=0\)
would be solvable, \emph{which is absurd}. --- Cases 3 and 6 of
Article 131 follow from this.

V.\ Let \(p,q\) be primes of the form \(4n+3\). We cannot have
\(p\,N\,q,q\,N\,p\) simultaneously. For suppose this were possible, and take
an auxiliary prime \(r\) of the form \(4n+1\) that is a nonresidue
of both \(p,q\). We would have \(qr\,R\,p,pr\,R\,q\); further (by II), \(p\,N\,r,q\,N\,r\),
giving \(pq\,R\,r\) and \(-pq\,R\,r\). Thus the equation \(-pxx-qyy+rzz=0\) would be
possible, contrary to the preceding observation. Case 8 of Article
131 follows from this.
